\documentclass[11pt]{article}

\usepackage[margin=1in]{geometry}
\usepackage[T1]{fontenc}
\usepackage[utf8]{inputenc}
\usepackage{amsmath}
\usepackage{amssymb}
\usepackage{booktabs}
\usepackage{enumitem}
\usepackage{graphicx}
\usepackage{hyperref}
\usepackage[numbers,sort&compress]{natbib}
\usepackage{xcolor}
\usepackage{listings}


\lstset{
  basicstyle=\ttfamily\footnotesize,
  breaklines=true,
  frame=single,
  numbers=none,
  showstringspaces=false,
  tabsize=2,
  language=Python,
  keywordstyle=\color{blue},
  commentstyle=\color{gray},
  stringstyle=\color{red}
}

\hypersetup{
  colorlinks=true,
  linkcolor=blue,
  citecolor=blue,
  urlcolor=blue,
  pdftitle={ZoraASI: A Governed Prototype for Bounded AI-Agent Workflows},
  pdfauthor={Christopher Michael Baird}
}

\title{\textbf{ZoraASI: A Governed Prototype for Bounded AI-Agent Workflows}\\[6pt]
\large A Contribution to A Theory of Everything}
\author{Christopher Michael Baird\\ZoraOS Research Project}
\date{July 20, 2026}

\begin{document}
\maketitle

\begin{abstract}
This paper presents ZoraASI, the experimental governed operating profile of ZoraOS, a local-first and model-agnostic prototype for bounded AI-agent workflows. The implementation separates replaceable model providers from user-owned memory, tool interfaces, execution budgets, access controls, and audit records. It includes a bounded tool-use loop, per-task iteration/token/tool-call budgets, default denial of unclassified tools, an in-memory SHA-256 hash-chained audit ledger, and REST endpoints for task submission, cancellation, monitoring, and audit retrieval. A four-machine topology (M5 orchestrator, M4 inference, M2 automation, M1 memory) is documented as a deployment design, not as a verified live cluster. In one observed run, a 6,340-page dissertation corpus (Zenodo record 21313827) had been indexed into 6,607 semantic chunks; a semantic query retrieved two passages with reported scores above 0.79, and a subsequent research-agent run completed five tool-loop iterations. The test suite provides unit-level evidence for selected controls but does not establish production security, factual reliability, consciousness, general intelligence, or safe unattended autonomy.
\end{abstract}

\section{Introduction}

Large language models can produce useful plans, answer questions, and generate code. But a model invocation is not a system. A system requires persistent memory, controlled access to tools, bounded execution, observability, and user authority. ZoraOS is an architecture designed to make these capabilities durable while allowing the underlying model provider to change without rewriting agent logic or tool schemas.

ZoraASI is the name used for the governed operating profile of ZoraOS. It specifies engineering controls\ ---\ least privilege, bounded tasks, explicit approval, audit records, and artifact classification\ ---\ for supervised automation. The name is not an empirical claim of artificial superintelligence. The goal is a system that can reason over user-authorized memory and tools while remaining bounded, observable, and interruptible.

This paper documents the ZoraASI prototype as observed and tested on July 20, 2026. It makes the following contributions:

\begin{enumerate}[leftmargin=*]
  \item A modular architecture and four-machine deployment design for model-agnostic AI-agent orchestration.
  \item A governance framework\ ---\ \texttt{AuditLedger}, \texttt{ExecutionBudget}, \texttt{ExecutionContext}, \texttt{GovernancePolicy}\ ---\ that enforces safety constraints at every tool invocation.
  \item A bounded tool-use execution loop with per-task iteration, token, and tool-call budgets.
  \item A full REST API with SSE streaming, task cancellation, and audit retrieval.
  \item Integration with a 6,340-page technical dissertation corpus (A Theory of Everything) indexed for semantic retrieval.
  \item Publication infrastructure including Zenodo record references, GitHub public/private separation, and reproducible documentation.
\end{enumerate}

\section{System Architecture}

\subsection{Multi-Machine Topology}

ZoraOS is designed to support four physical machines, each serving a distinct role. The
topology in Table~\ref{tab:topology} is architectural; the reviewed session did not
verify live distributed orchestration across all four nodes.

\begin{table}[ht]
\centering
\caption{Deployment topology across four machines.}
\label{tab:topology}
\begin{tabular}{p{0.12\linewidth} p{0.22\linewidth} p{0.56\linewidth}}
\toprule
\textbf{Node} & \textbf{Primary Role} & \textbf{Services} \\
\midrule
M5 & Orchestrator & FastAPI gateway, planner, agent manager, UI dashboard, SSE streaming \\
M4 & Inference & Ollama, embedding models, inference workloads, OCR processing \\
M2 & Automation & Authorized research jobs, build pipelines, version control operations \\
M1 & Memory & PostgreSQL, Redis, vector store, document persistence, knowledge graph \\
\bottomrule
\end{tabular}
\end{table}

This separation is an operational design choice, not a requirement for minimal deployment. The architecture supports both local inference services and hosted providers through a common \texttt{ModelManager} interface.

\subsection{Core Data Flow}

The system operates on a structured task model:

\begin{enumerate}[leftmargin=*]
  \item A user submits a task through the dashboard or REST API, specifying agent type, goal, authorized tools, and optional budget.
  \item The gateway applies loopback/API-key access control, creates a governed \texttt{ExecutionContext}, and records a \texttt{task\_started} audit event.
  \item The planner decomposes the goal into steps.
  \item The router selects an optimal model and provider based on the task type.
  \item An agent receives the system prompt, task objective, model configuration, and authorized tool set.
  \item The agent alternates between model responses and governed tool results until completion or a bound is reached.
  \item Results, metrics, and audit events are returned. In the active prototype, task events are recorded in a per-process in-memory hash chain.
\end{enumerate}

\section{Governance Framework}

The governance layer is the central engineering contribution of ZoraASI. It implements a default-deny security model with capability classification, per-task budgets, and hash-chained audit provenance.

\subsection{Capability Classification}

Every registered tool is classified into one of five code-level capability classes. The
table separates proposed operational sub-classes for clarity; unknown tools map to
\texttt{PROHIBITED} and are denied.

\begin{table}[ht]
\centering
\caption{Tool capability classification and required controls.}
\label{tab:classes}
\begin{tabular}{p{0.22\linewidth} p{0.38\linewidth} p{0.30\linewidth}}
\toprule
\textbf{Class} & \textbf{Examples} & \textbf{Required Control} \\
\midrule
Read-only & memory\_search, pdf\_reader & Task scope, audit event \\
Local reversible write & filesystem, memory\_write, git & Explicit task approval \\
Local irreversible write & data deletion, config modification & Explicit confirmation, backup \\
External communication & web\_search, web\_post & Per-action confirmation \\
Third-party control & eq\_send\_keys, eq\_read\_screen & Sandbox only, session supervision \\
Prohibited & packet inspection, evasion, credential harvesting & Denied unconditionally \\
\bottomrule
\end{tabular}
\end{table}

The \texttt{GovernancePolicy} class implements this classification and the authorization logic:

\begin{lstlisting}
class CapabilityClass(StrEnum):
    READ_ONLY = "read_only"
    LOCAL_WRITE = "local_write"
    EXTERNAL = "external"
    THIRD_PARTY_CONTROL = "third_party_control"
    PROHIBITED = "prohibited"

class GovernancePolicy:
    _classes: dict[str, CapabilityClass] = {
        "memory_search": CapabilityClass.READ_ONLY,
        "pdf_reader": CapabilityClass.READ_ONLY,
        "filesystem": CapabilityClass.LOCAL_WRITE,
        "web_search": CapabilityClass.EXTERNAL,
        "eq_send_keys": CapabilityClass.THIRD_PARTY_CONTROL,
    }

    def classify(self, tool_name: str) -> CapabilityClass:
        return self._classes.get(tool_name, CapabilityClass.PROHIBITED)
\end{lstlisting}

Authorization follows a strict chain: prohibited and unclassified tools are denied; if
no task context exists, only registered read-only tools are permitted; if the task is
cancelled, all tools are denied; third-party control tools are disabled by the active
gateway policy; and local-write or external tools require explicit per-task approval.

\subsection{Execution Budget}

Every task carries an \texttt{ExecutionBudget} that imposes limits on resource consumption:

\begin{lstlisting}
@dataclass
class ExecutionBudget:
    max_tool_calls: int = 20
    max_iterations: int = 10
    max_tokens: int | None = None
    tool_calls_used: int = 0

    def consume_tool_call(self) -> bool:
        if self.tool_calls_used >= self.max_tool_calls:
            return False
        self.tool_calls_used += 1
        return True
\end{lstlisting}

\subsection{Execution Context}

The \texttt{ExecutionContext} binds the budget, approved tools, and cancellation state to a specific task:

\begin{lstlisting}
@dataclass
class ExecutionContext:
    task_id: str
    approved_tools: frozenset[str]
    budget: ExecutionBudget
    cancelled: bool = False

    def permits(self, tool_name: str) -> bool:
        return tool_name in self.approved_tools
\end{lstlisting}

\subsection{Audit Ledger}

The active \texttt{AuditLedger} provides an in-memory, SHA-256 hash-chained event store.
Each event cryptographically links to its predecessor, enabling integrity checks while
the process remains alive:

\begin{lstlisting}
@dataclass(frozen=True)
class AuditEvent:
    id: str
    timestamp: float
    task_id: str | None
    event_type: str
    payload_digest: str
    previous_digest: str
    digest: str
\end{lstlisting}

The digest is computed as:

\begin{equation}
  d_i = \mathrm{SHA256}(e_i : t_i : \mathrm{task\_id}_i : \mathrm{type}_i : \mathrm{payload\_digest}_i : d_{i-1})
\end{equation}

where \(d_0 = 0^{64}\). Verification replays the chain and confirms that every event's
previous\_digest matches the computed digest of its predecessor and that every event's
own digest is recomputable. This is tamper-evident within the process, not durable
storage. An experimental PostgreSQL implementation serializes writers with an advisory
lock and verifies rows in Python, but it is not yet wired into the gateway.

\subsection{Governed Tool Execution}

The \texttt{ToolManager} integrates policy, context, budget, and audit into a single execution path. Every tool call passes through this pipeline:

\begin{lstlisting}
async def execute(self, tool_name: str, **kwargs) -> ToolResult:
    context = self._execution_context.get()
    if self._policy:
        allowed, reason = self._policy.authorize(tool_name, context)
        if not allowed:
            self._audit_ledger.record("tool_denied",
                {"tool": tool_name, "reason": reason}, task_id)
            return ToolResult(success=False, error=reason)
    if context and not context.budget.consume_tool_call():
        self._audit_ledger.record("budget_exceeded", ...)
        return ToolResult(success=False, error="budget exceeded")
    # execute tool, record tool_started / tool_completed / tool_failed
\end{lstlisting}

\section{Bounded Tool-Use Loop}

For a task objective \(g\), a system prompt \(s\), and authorized tools \(T\), the agent initializes a conversation with \((s, g)\). At iteration \(i\), the model yields either a final response or a set of tool calls. Each tool result is serialized into the conversation before the next model invocation. The loop terminates on a final response or at a configured iteration limit \(I_{\max}\):

\begin{equation}
  \mathrm{terminate} \iff \mathrm{final\_response} \lor i \geq I_{\max}
\end{equation}

Token budget enforcement is also supported. If the cumulative token count exceeds a per-task limit, the agent terminates with an explicit budget-exceeded error.

The structured result includes:
\begin{itemize}[leftmargin=*]
  \item A success flag
  \item Task identifier and agent name
  \item Model identifier (when supplied by the provider)
  \item Iteration count and token usage (when reported)
  \item Wall-clock latency in milliseconds
  \item Full tool-call trace (tool name, arguments, success, output)
\end{itemize}

This structure makes failure states observable rather than silently extending a conversation indefinitely.

\section{API Layer}

The REST API exposes the full system programmatically:

\begin{table}[ht]
\centering
\caption{REST API endpoints.}
\label{tab:api}
\begin{tabular}{p{0.28\linewidth} p{0.22\linewidth} p{0.40\linewidth}}
\toprule
\textbf{Endpoint} & \textbf{Method} & \textbf{Purpose} \\
\midrule
\texttt{/agents/run} & POST & Submit agent task with approved tools and budget \\
\texttt{/agents/run/stream} & POST & SSE-streaming agent execution \\
\texttt{/agents/tasks} & GET & List all tasks \\
\texttt{/agents/tasks/\{id\}} & GET & Get task status and result \\
\texttt{/agents/tasks/\{id\}/cancel} & POST & Cancel a running task \\
\texttt{/agents/tasks/\{id\}/audit} & GET & Retrieve audit events for a task \\
\bottomrule
\end{tabular}
\end{table}

The SSE streaming endpoint emits events for task start, plan creation, route selection, tool calls, and task completion\ ---\ providing real-time visibility into agent execution.

Task cancellation sets a cancellation flag in the \texttt{ExecutionContext} that prevents all subsequent tool invocations, and records a \texttt{task\_cancel\_requested} audit event.

\section{Memory and Corpus Integration}

The semantic memory subsystem indexes documents into collections for retrieval-augmented generation. A significant case study was performed using the dissertation \emph{A Theory of Everything} (Zenodo record 21313827), a 6,340-page technical work by C. M. Baird.

The reviewed local deployment reported the following observations:
\begin{itemize}[leftmargin=*]
  \item The dissertation had been indexed into 6,607 chunks in a \texttt{physics} collection.
  \item A semantic query for ``quantum gravity'' returned two ranked passages with reported scores 0.8061 and 0.7911.
  \item A subsequent research-agent task completed five tool-loop iterations and generated a synthesis after memory search/read operations.
\end{itemize}

This is a single-deployment functional observation, not an independent scientific
validation of the corpus, a retrieval benchmark, or evidence of access to the full body
of human knowledge. Reproducible evaluation requires a frozen corpus manifest, query
set, relevance judgments, and independent reruns.

\section{Testing and Verification}

At import, the repository's baseline suite reported 35 passing tests and six
\texttt{datetime.utcnow()} deprecation warnings. The reviewed tests covered:

\begin{itemize}[leftmargin=*]
  \item \textbf{Governance enforcement}: external tools require explicit approval; denied tools produce audited error responses.
  \item \textbf{Third-party control restriction}: desktop-control tools (e.g., \texttt{eq\_send\_keys}) are denied even when listed in approved tools, unless an approved sandbox is present.
  \item \textbf{Budget enforcement}: tasks that exceed their tool-call budget are terminated with a budget-exceeded error.
  \item \textbf{Audit chain integrity}: all test scenarios verify that the \texttt{AuditLedger} passes \texttt{verify()}, confirming hash-chain integrity.
  \item \textbf{Memory operations}: document creation, default collection assignment, and search result correctness.
  \item \textbf{Router fallback}: provider selection degrades gracefully when the configured default is unavailable.
\end{itemize}

These results are unit-level evidence from one local environment. They do not establish
production security, cross-platform behavior, live provider reliability, or
multi-machine correctness.

\subsection{Governance Test Results}

The governance test suite (\texttt{tests/test\_governance.py}) demonstrates three critical guarantees:

\begin{enumerate}[leftmargin=*]
  \item \textbf{Approval requirement (test\_external\_tool\_requires\_explicit\_approval)}: an external tool (\texttt{web\_search}) is denied without a governed task context. With a context that explicitly approves the tool, it executes successfully. The audit ledger records \texttt{tool\_started} and \texttt{tool\_completed} events, and the chain verifies.
  \item \textbf{Third-party denial (test\_third\_party\_desktop\_control\_is\_denied\_even\_when\_approved)}: even when \texttt{eq\_send\_keys} is listed in the approved tools set, the active policy denies it. Gateway support for a dedicated approved sandbox has not been implemented.
  \item \textbf{Budget enforcement (test\_tool\_call\_budget\_is\_enforced)}: a task with \texttt{max\_tool\_calls=1} executes the first tool call successfully; the second is denied with a budget-exceeded error. The audit ledger chain verifies.
\end{enumerate}

\section{Publication and Artifact Infrastructure}

ZoraASI is accompanied by a publication artifact structure intended to support reproducibility and provenance:

\begin{table}[ht]
\centering
\caption{Publication and artifact infrastructure.}
\label{tab:pub}
\begin{tabular}{p{0.26\linewidth} p{0.64\linewidth}}
\toprule
\textbf{Artifact} & \textbf{Purpose} \\
\midrule
Technical report & \texttt{paper/main.tex} and \texttt{references.bib} \\
ZoraASI paper & \texttt{paper/zoraasi\_complete.tex} \\
Product requirements & \texttt{docs/pdr/} sources and rendered artifacts \\
Citation metadata & \texttt{CITATION.cff} \\
Security policy & \texttt{SECURITY.md} \\
Governance baseline & \texttt{docs/GOVERNANCE.md} \\
Artifact rules & \texttt{docs/ARTIFACTS.md} \\
Publication index & \texttt{docs/PUBLICATION.md} \\
Zenodo references & \texttt{docs/ZENODO\_21313827.md} and \texttt{docs/ZENODO\_21464562.md} \\
\bottomrule
\end{tabular}
\end{table}

\subsection{Theory of Everything Corpus Reference}

The corpus used for the semantic retrieval case study is the dissertation \emph{A Theory of Everything} by C. M. Baird \citep{baird2026dissertation_zenodo}. This 6,340-page work is hosted on Zenodo (record 21313827). The ZoraASI OS publication \citep{baird2026zoraasi_zenodo} (Zenodo record 21464562) is the companion primary publication documenting the governed AI operating system architecture. Local PDF copies of both are included in the ZoraOS repository under \texttt{docs/paper/} for citation and provenance context. The authoritative records are at:

\begin{itemize}
  \item Dissertation: \url{https://zenodo.org/records/21313827}
  \item ZoraASI OS: \url{https://zenodo.org/records/21464562}
\end{itemize}

\subsection{GitHub Repository Separation}

The codebase is split into two repositories:
\begin{itemize}[leftmargin=*]
  \item \textbf{Public} (\texttt{github.com/Cbaird26/zoraos}): source code, documentation, tests, and publication artifacts\ ---\ all free of credentials and private data.
  \item \textbf{Private evidence} (\texttt{github.com/Cbaird26/zoraos-private-evidence}): protocol templates and artifact manifests for experimental evidence (no game logs, screenshots, or personally identifying information).
\end{itemize}

.gitignore rules exclude \texttt{.env} files, runtime data, logs, screenshots, node modules, and generated PDFs from version control, ensuring that only reviewed, classified artifacts enter the public record.

\section{ZoraASI Safety Principles}

The governance profile rests on seven operating principles \citep{shneiderman2020human, nist2023airmf}:

\begin{enumerate}[leftmargin=*]
  \item \textbf{Human authority}: the operator sets goals, permissions, budgets, retention rules, and stop conditions.
  \item \textbf{Least privilege}: tools begin unavailable and are granted only the minimum scope needed for a task.
  \item \textbf{Bounded execution}: every task has limits for iterations, wall-clock time, model cost, and tool calls.
  \item \textbf{Observable actions}: inputs, tool calls, approvals, outputs, failures, and stop events require durable records.
  \item \textbf{Reversible defaults}: read-only analysis and proposals are preferred over external side effects.
  \item \textbf{User-owned memory}: the user controls memory ingestion, retrieval, export, retention, and deletion.
  \item \textbf{No policy evasion}: the system must not bypass technical controls, platform policies, access restrictions, or consent requirements.
\end{enumerate}

\subsection{Prohibited Mechanisms}

The following are unconditionally prohibited:
\begin{itemize}[leftmargin=*]
  \item Network packet interception or injection
  \item Reading or writing application memory outside the ZoraOS process
  \item Process injection or code injection
  \item Anti-cheat evasion or platform control bypass
  \item Credential extraction or harvesting
  \item Unattended activity on third-party services
  \item Escalation of privilege beyond user's authority
\end{itemize}

\subsection{Implementation Status}

\begin{table}[ht]
\centering
\caption{Implemented and required governance controls.}
\label{tab:controls}
\begin{tabular}{p{0.36\linewidth}p{0.23\linewidth}p{0.30\linewidth}}
\toprule
\textbf{Control} & \textbf{Current Status} & \textbf{Required Next Step} \\
\midrule
Iteration/token/tool-call budgets & Prototype & Add validated cost and wall-clock enforcement \\
Registered tool selection & Implemented & Add permission scopes and expiry \\
Local kill-switch convention & Prototype & Add visible control, integration test, and terminal task state \\
Tool-call reporting and in-memory hash chain & Prototype & Integrate durable storage and retention \\
Desktop screen capture and OCR & Prototype & Evaluate only in an authorized sandbox \\
External action confirmation & Proposed & Implement a mandatory approval broker \\
PostgreSQL audit module & Experimental, not integrated & Add live database integration tests and gateway wiring \\
Permission broker & Proposed & Add capability scopes, actor identity, and expirations \\
\bottomrule
\end{tabular}
\end{table}

\section{Discussion}

\subsection{Architectural Significance}

ZoraASI represents a departure from the dominant paradigm of AI deployment. Rather than treating the model as the product, it treats the governed operating layer as the enduring asset. This has important consequences:

\begin{itemize}[leftmargin=*]
  \item \textbf{Provider abstraction}: providers can be swapped without rewriting agent logic, although live parity across providers has not been established.
  \item \textbf{User sovereignty}: memory, tools, and policies are owned and controlled by the user, not by a cloud provider.
  \item \textbf{Reproducibility direction}: the system produces structured results and audit digests; independent verification still requires frozen inputs, durable records, and rerun instructions.
  \item \textbf{Bounded-failure direction}: budgets, cancellation flags, and default denial reduce some failure modes but do not constitute a complete safety guarantee.
\end{itemize}

\subsection{Limitations}

The system has important limitations:

\begin{itemize}[leftmargin=*]
  \item The active audit ledger is in-memory and does not survive process restart. The experimental PostgreSQL module is not integrated into the gateway; durable storage, migration tests, export, retention, and independent verification remain required.
  \item The permission broker lacks persistence: approved tool scopes and expirations are not yet stored or enforced across sessions.
  \item Desktop interaction research included keystrokes sent to a live third-party commercial application. Generic OCR did not reliably extract text. This was not a controlled sandbox or a valid safety evaluation and must not be repeated as unattended or detection-evasion testing.
  \item The imported API used a placeholder gateway secret and was bound to all interfaces. The revised prototype restricts placeholder-key access to loopback and requires an API key when configured, but production identity, authorization, rate limiting, and transport security remain absent.
  \item All tests are unit-level. Integration tests with live model providers and multi-machine deployment tests have not been performed.
  \item The system has not been evaluated for long-term unattended operation. The governance profile explicitly disclaims reliable autonomous competence.
\end{itemize}

\subsection{Relationship to the Theory of Everything}

The semantic memory case study indexed the full 6,340-page \emph{A Theory of Everything} dissertation and successfully retrieved relevant passages in response to queries. This demonstrates that ZoraOS can serve as a computational substrate for working with the complete corpus\ ---\ enabling search, synthesis, and iterative reasoning over the text.

The work described in this paper is not itself a contribution to physics or unified field theory. It is an infrastructure contribution: a governed, reproducible, user-sovereign platform upon which such work can be conducted safely and transparently. The Theory of Everything provides the content; ZoraASI provides the container.

\section{Conclusion}

ZoraASI is a documented and unit-tested prototype for governed AI-agent workflows. It implements:

\begin{itemize}[leftmargin=*]
  \item A proposed multi-machine architecture separating orchestration, inference, automation, and memory.
  \item A default-deny governance framework with capability classification, per-task budgets, and SHA-256 hash-chained audit.
  \item A bounded tool-use execution loop with structured results and failure observability.
  \item A full REST API with streaming, cancellation, and audit retrieval.
  \item Integration with a 6,340-page technical corpus for retrieval-augmented generation.
  \item A local automated test suite covering selected governance and core behaviors.
  \item Publication artifacts including Zenodo references, CITATION.cff, and public/private GitHub separation.
\end{itemize}

The system is not a claim of artificial general intelligence, consciousness, or safe unattended autonomy. It is an engineering prototype and research agenda for bounded, user-controlled AI operation. The observed results justify further controlled evaluation of provider-independent orchestration, retrieval, and policy enforcement; they do not demonstrate that safe persistent autonomy has been achieved.

The source code and documentation are available at \url{https://github.com/Cbaird26/zoraos}. The dissertation corpus that grounds the system's knowledge is at \url{https://zenodo.org/records/21313827}. The primary ZoraASI OS publication is at \url{https://zenodo.org/records/21464562}.

\bibliographystyle{plainnat}
\bibliography{references}

\end{document}
